\section{Test Infrastructure}
\label{sec:testing}

Beyond the 100 KAT vectors from the C reference implementation (all of
which verify successfully), we maintain a growing suite of negative and
vulnerability test vectors in the C2SP/wycheproof JSON
format~\cite{wycheproof}.  These vectors are designed to detect not
just bugs but \emph{vulnerabilities}---they exercise rejection paths
that the KAT vectors, being all valid, never touch.

Invalid vectors are independently generated by byte-level perturbation
(no dependency on the C reference implementation).  Valid vectors cite
the C reference KAT file as their provenance; we intend to replace
them with independently Sage-generated vectors for eventual upstream
contribution to the C2SP/wycheproof project.

\subsection{Algorithm naming}

We identify the parameter set as \texttt{SQIsign\_248}, where 248 is
the torsion exponent~$f$ in $p = 5 \cdot 2^{248} - 1$.  The exponent
is the fundamental parameter from which all others derive: field size,
isogeny degrees, key sizes, signature sizes.  Unlike ``NIST-I'' (a
security-level label that could be reassigned to different parameters)
or ``353'' (a derived quantity---the secret key size in bytes---that
would change if the encoding changes), $f = 248$ pins the actual
cryptographic parameters.

\subsection{Verification vectors}

Table~\ref{tab:verify-vectors} summarizes the verification vector
categories.  Each vector consists of a public key, message, signature,
and expected result (\texttt{valid} or \texttt{invalid}).  Invalid
vectors are derived from valid KAT vectors by targeted perturbation.

\begin{table}[ht]
\centering
\small
\begin{tabular}{@{}llp{0.46\textwidth}@{}}
\toprule
\textbf{Category} & \textbf{Count} & \textbf{What it detects} \\
\midrule
Normal (valid KAT) & 2 & Baseline: C~ref signatures verify. \\
WrongMessage & 2 & Different, empty, extended, and truncated
  messages. \\
WrongPublicKey & 1 & Cross-key verification. \\
BitFlip (per field) & 7 & Single-bit corruption in each of the seven
  signature fields. \\
AllZero / AllOnes & 2 & Degenerate byte patterns. \\
ExponentOverflow & 4 & $n_{\mathrm{bt}} + r_{\mathrm{rsp}} > e_{\mathrm{rsp}}$,
  causing $e'_{\mathrm{rsp}}$ underflow. \\
NonCanonicalFp & 3 & Field elements $\ge p$ in \texttt{E\_aux} or pk.
  Silently reduces mod~$p$, producing the wrong curve. \\
SingularCurve & 3 & $A = \pm 2$ (singular Montgomery curve,
  discriminant zero) in \texttt{E\_aux} or pk. \\
SpecialJInvariant & 1 & $A = 0$ ($j = 1728$, extra automorphisms
  $\pm 1, \pm i$). \\
StartingCurve & 1 & \texttt{E\_aux} = $E_0$ ($A = 6$). \\
DegenerateMatrix & 2 & \texttt{M\_chl} all-zero or identity. \\
ZeroChallenge & 1 & Challenge $= 0$ (degenerate Fiat--Shamir). \\
ScalarOverflow & 1 & Matrix entry $\ge 2^{e_{\mathrm{rsp}}}$. \\
ChallengeOverflow & 1 & High bits set beyond $e_{\mathrm{chl}} = 122$. \\
SplicedField & 3 & Signature field replaced with the corresponding
  field from a different valid signature. \\
MalformedPk & 2 & All-zero or all-\texttt{0xFF} public key. \\
\midrule
\textbf{Total} & \textbf{39} & \\
\bottomrule
\end{tabular}
\caption{Verification test vector categories for
  \texttt{SQIsign\_248} (C2SP/wycheproof JSON format).}
\label{tab:verify-vectors}
\end{table}

\paragraph{Finding: verification panics on malformed input.}
Several invalid vectors (bit flips in \texttt{E\_aux},
\texttt{M\_chl}, and \texttt{chl}) reach the $(2,2)$-isogeny
splitting step (Algorithm~8.42, \textsc{GetIndexSplitting}), which
uses a \texttt{debug\_assert!} to check that exactly one null-point
index is zero.  Malformed signatures violate this invariant, and the
assertion panics rather than returning an error.  In debug builds this
is a denial-of-service vector: any untrusted signature can crash the
verifier.  In release builds the assertion compiles out and the
function returns a wrong index silently.  The correct fix is to return
\texttt{Option<usize>} and propagate the error as
\texttt{VerificationFailed}.  Our test runner uses
\texttt{std::panic::catch\_unwind} as a temporary workaround.

\subsection{Key generation vectors}

We maintain 12 keygen/sk-parsing vectors:
\begin{itemize}[nosep]
  \item Deterministic keygen from KAT seeds produces the expected
    (pk, sk) byte-for-byte.
  \item The pk embedded in the first 65 bytes of sk matches the
    standalone pk.
  \item Malformed sk rejection: truncated, extended, all-zero,
    all-\texttt{0xFF}, single-bit flips in each region (pk, ideal
    norm, M\textsubscript{sk} matrix), and a ``Frankenstein'' sk
    with the pk from one vector spliced onto the sk body of another.
\end{itemize}

\subsection{Signing vectors (planned)}

The signing vector file is a skeleton with one vector (C~ref KAT
round-trip) and a roadmap for categories to add as signing stabilizes:
\begin{itemize}[nosep]
  \item \textbf{NonceIndependence}: same (sk, msg) with different
    randomness must produce distinct commitment curves.  Nonce reuse
    leaks the secret key.
  \item \textbf{CornacchiaBranching}: inputs crafted so that
    \textsc{RepresentInteger}'s Cornacchia subroutine takes different
    branches, targeting SPA
    vulnerabilities~\cite{EPRINT:MCJKR25}.
  \item \textbf{ResponseEdgeCases}: signatures with $r_{\mathrm{rsp}}
    > 0$ or $n_{\mathrm{bt}} > 0$, exercising the even-response
    isogeny and backtracking paths that are rare in KAT vectors.
\end{itemize}

\subsection{Mutation testing}

We use \texttt{cargo-mutants}~\cite{cargo-mutants} to verify that
the test suite catches bugs, not just that it runs.  Mutation testing
inserts faults (replacing operators, return values, conditionals) and
checks whether any test fails.  A surviving mutant is a missing test.

Running \texttt{cargo-mutants} on the \texttt{keys/} module (212
mutants) found 8 surviving mutants---6 in \texttt{ChallengeMatrix::from\_bytes}
validation and 2 in keygen---all of which represented real test gaps.
Writing targeted tests killed the validation mutants; the keygen
mutants survive because keygen tests are slow (each invocation runs
a $(2,2)$-isogeny chain) and are marked \texttt{\#[ignore]}.

CI runs incremental mutation testing on every pull request (only
changed code via \texttt{-{}-in-diff}) and a full sharded run weekly.
